\begin{center}
{\large\bfseries Abstract}
\end{center}

\noindent We set out to find which personal and regional factors go together with lung cancer in US adults, with an answer we could defend line by line. We stacked five years of the National Health Interview Survey (NHIS 2019--2023; \nAdults adults after cleaning) and joined three regional sources: EPA fine-particle pollution (PM2.5), CDC adult smoking prevalence (BRFSS) and official lung-cancer incidence (US Cancer Statistics). Lung cancer is rare in these data -- \nCases adults, \sampleRate\% of the sample and 0.31\% after survey weighting -- and that rarity shaped almost every design decision, starting with rejecting accuracy as a metric (an ``always no'' rule would already be 99.59\% accurate). A detailed profile of the study population (Section~\ref{sec:profile}) shows who the cases are: 76\% are aged 65 or older, 82\% have ever smoked and 47\% have COPD, against 31\%, 38\% and 5\% of the other adults.

Age, smoking history and COPD stand out clearly. Ever-smokers have \orSmoking times the odds of never-smokers (95\% CI \orSmokingLo--\orSmokingHi), and still \orSmokingMH times after adjusting for age group; people with lung cancer are much older (median 72 vs 54 years; a random case is older than a random non-case 79.5\% of the time). COPD carries an odds ratio of \orCOPD; among adults with known smoking status the crude \orCOPDcrude falls to \orCOPDsmoking after adjusting for smoking, and among those whose age is also known it falls to \orCOPDsmokingAge after adjusting for smoking and age group together. Lower education and income, underweight and poor self-rated health also go with higher rates, the last two most likely through reverse causation. Of \nTests hypothesis tests, \nSignificant remain significant after Holm correction, and a survey-weighted, cluster-robust check leaves every key odds ratio far from 1. At the regional level, the survey rate agrees with the official registry rate (Spearman $\rho = \rhoRegistry$ over 20 region-years), but this agreement comes from differences between the four regions rather than from year-to-year changes, so it is consistent with -- not proof of -- a valid target; PM2.5 and regional smoking prevalence show no detectable association, and as features the three regional columns carry a single shared signal (which region-year a person lives in) rather than three separate ones.

The result is a reproducible pipeline that one command rebuilds in under two minutes, with every cleaning step logged, all learned steps fitted on training rows only (checked by an automated assertion), and a leakage-free set of \nFeatures features selected from \nCandidates candidates. Feature selection also removes dependencies that span several columns -- the comorbidity count, which is the exact sum of five condition flags, one reference category per one-hot group, and the region dummies -- so that no feature has a variance inflation factor above \maxVIFshort and a linear model can be fitted without perfect multicollinearity. A principal component analysis confirmed that the information is spread over many dimensions (\pcaNinetyFive of \nFeatures components for 95\% of the variance, and no near-zero component), so we keep the original, interpretable features for the modelling phase.

\medskip
\noindent\textbf{Keywords:} lung cancer, risk factors, NHIS, CRISP-DM, data integration, hypothesis testing, confounding, class imbalance, data leakage, PCA

\clearpage
